\documentclass[12pt]{article}

\usepackage[margin=1in]{geometry}
\usepackage{setspace}
\usepackage{natbib}
\usepackage{url}
\setstretch{1}

\setlength{\parindent}{0pt}
\setlength{\parskip}{0.8em}

\title{Per-Launch Efficiency or System-Level Sustainability: Evaluating Environmental Strategies for LEO Broadband Expansion}
\author{Wenqin Huang}
\date{\today}

\begin{document}

\maketitle

%TC:ignore
\begin{center}
\small \textbf{Word Count: 2,132}
\end{center}
%TC:endignore


In today's era, LEO broadband constellations are not simply symbolic space activities, but important components of communication infrastructure. The 2026 stimulus material from the U.S. Surgeon General's advisory explains that social connection influences individual health through biological, psychological, and behavioral pathways, and that stronger social networks are linked with better mental health, cognitive function, education, and economic outcomes~\citep{us_surgeon_general_2023_social_connection_health}. This stimulus does not directly discuss satellites, but it provides a link between communication access and public health significance. Under this context, another 2026 stimulus material, Osoro et al.'s research on LEO broadband megaconstellation, connects this social value to large scale satellite deployments: in remote areas where traditional internet service does not reach, LEO satellites can provide internet connections, improve social connections, and subsequently support mental health~\citep{osoro2025leo,sandau2010smallsat}. Therefore, the rapid expansion of OneWeb and Starlink is motivated partly by the demand of globally-covered connections. Thus, space launches should not simply be stopped and cancelled.



However, the same process that expands global connection also increases environmental and orbital costs. Space activity is no longer occasional; it is becoming a rapidly expanding industrial sector: according to Ryan et al.'s 2022 research, its scale will leap from 350 billion in 2019 to over a trillion by 2040~\citep{ryan2022ozone,bushnell2018newspace}. This expansion will probably cause a substantial increase in its environmental impacts. Additionally, through rocket launching activities, multiple pollutants will be directly injected into the upper atmosphere, where pollution may contribute to ozone destructions~\citep{ryan2022ozone,kokkinakis2022pollution}. Moreover, this expansion is not only something that will happen in the future, but is already happening: the frequency of rocket launches increased considerably from 58 launches in 2003 to 329 launches in 2025~\citep{mcdowell2025space}. The environmental impact is not limited to Earth's atmosphere. Satellite constellations will also occupy Earth's orbits, create a large amout of space debris, and subsequently obstruct future satellite deployments~\citep{torres2020reusable}. Thus, the two 2026 stimulus sources and the space-launch evidence together create the central tension of this research: megaconstellations can strengthen social connection, but the launches and orbital occupation behind them may also create considerable sustainability costs.



Based on this tension, the true essential problem is ``to what extent can technological strategies reduce these environmental costs and improve sustainability'', rather than ``whether people should continue launching or not''. To address the environmental costs caused by frequent space launches, various technological strategies have been proposed to improve the sustainability of launch systems. Among these, reusable launch vehicles, alternative propellants, and satellite miniaturization are often considered to be most practical approaches. Their effects can be evaluated not only by whether they reduce per-launch or per-satellite impacts, but also by whether these reductions still remain meaningful when LEO satellite deployment continues expanding~\citep{wilson2022lifecycle,osoro2025leo}. This paper argues that, although these technological strategies can reduce environmental impacts at the per-unit level, their ability to create system-wide sustainability remains limited by financial incentives, adoption barriers, industrial inertia, and rebound effects from LEO megaconstellation expansion.



Using alternative propellants is the most direct technique under an emissions-based criteria, since they change the combustion outputs of launch vehicles directly, rather than only improving the design of the vehicle. Dallas et al. first construct this perspective by reviewing the environmental impact generated by traditional propellants: solid fuel, which typically involves chemicals like ammonium perchlorate, may potentially damage the ozone the most, while kerosene-based propellants mainly emit carbon dioxide (CO$_2$) and black carbon, which may intensify the Greenhouse effect~\citep{dallas2020review}. Barato then supports this comparison by suggesting liquid hydrogen as a future propulsion candidate, since its emissions mainly contain only water vapor and have negligible ecological toxicity~\citep{barato2023fuels}. Hai et al. also extend this argument by suggesting that environmentally friendly propellants not only reduce toxic products, but also lower the risks of transportation and operation~\citep{hai2026propellants}. To be more specific, Calabuig et al.'s life-cycle study gives this emissions view more numerical support: the carbon footprint of liquid hydrogen fleets is typically 2--8 times lower than that of methane-based vehicle fleets, and propellant choice is the main differentiation factor under environmental criteria~\citep{calabuig2024reusable}. This group of evidence is useful to the extent of direct combustion emissions, because it compares the chemical products and fleet-level fuel choices of different propellants. However, it is still less sufficient for evaluating system-wide, and overall sustainability, because proving that a cleaner propellant exists does not prove that the industry can adopt it at scale.

However, this direct-emissions view may still be complicated by life-cycle and implementation evidence. Dallas et al. argue that liquid hydrogen is inherently difficult to store: its density is lower, which means it requires a storage system that is significantly larger and more complex than traditional fuels; additionally, the boiling point of liquid hydrogen is considerably lower, so the propellant is more likely to vaporize and existing storage tanks may not be suitable~\citep{dallas2020review}. This indicates that the cleaner propellant may not be easier to store or use. Calabuig et al. also move the discussion from chemical performance to adoption feasibility, arguing that although sustainable propellants are environmentally beneficial, their cost is still significantly higher than that of traditional propellants, and the industry around production, transportation, and storage is still immature~\citep{calabuig2024reusable}. Since more rocket launching activities are now carried out by private companies, the higher cost may discourage companies from utilizing alternative propellants. In conclusion, alternative propellants offer the strongest per-launch emissions reduction, but their actual effect depends on the speed of industrial transformation and financial incentives, rather than only the environmental performance of the fuel itself.



However, propellant choice alone does not fully determine the environmental impact of space launches. When the evaluation moves from fuel choice to launch-system design, reusable launch vehicles offer a plausible solution because they reduce manufacturing emissions and material waste. Calabuig et al. argue that applying reusable rocket technology can reduce the manufacturing cost of launching, but it will not necessarily reduce the net environmental impact~\citep{calabuig2024reusable}. On one hand, reusable vehicles can improve the material efficiency of fixed equipment, and when evaluated under the criterion of launching impact per unit of mass, reusable vehicles can improve the efficiency both environmentally and financially. However, Calabuig et al. also limit this positive perspective with numerical evidence: for liquid hydrogen fleets, the climate impact could not be reduced by more than 40\%; for methane fleets, the impact is even less significant, around 20\%~\citep{calabuig2024reusable}. Wilson et al. further complicate the evaluation by arguing that long-term life-cycle assessment is critical in evaluation of the environmental impact of space missions~\citep{wilson2022lifecycle}. In this case, the question is not only whether one reusable launch is cleaner, but whether cheaper and more efficient launch systems encourage a greater total number of launches.

This concern of rebound effect can be supported by launch frequency data, although it should not be treated as direct proof. The FAA (Federal Aviation Administration) environmental reviews show that the annual Falcon launch cadence considered for SpaceX has risen steadily, from 20 launches per year at LC-39A and 50 at SLC-40 in the 2020 assessment, to 36 launches per year at Vandenberg in 2023, 50 in 2024, and 100 in the 2025 Vandenberg EIS~\citep{faa2025falconprogram}. McDowell's launch data also shows that SpaceX reached 165 launches in 2025 alone~\citep{mcdowell2025space}. These statistics do not prove that reusability alone caused the increase in launch frequency, but they illustrate why lower-cost reusable systems cause a serious rebound-effect. To summarize, reusable vehicles are environmentally useful and cost effective when measuring per launch or per kilogram depolyed, but their system-wide benefit depends on whether cost reductions lead to additional launches.



Satellite miniaturization improves the functional efficiency of individual satellites, but it also changes what should be evaluated from one satellite to the whole constellation. Sandau et al. first provide the access and efficiency side of this perspective, suggesting that satellite miniaturization could improve the efficiency per unit of mass by reducing mass while the function remains unchanged, and enable satellite broadband systems to cover globally~\citep{sandau2010smallsat}. This argument is useful when the goal is coverage, because smaller satellites make broadband deployment cheaper, easier to design, and easier to launch. It also connects with Osoro et al.'s argument that LEO broadband can improve access in remote and underserved regions where terrestrial internet service does not reach~\citep{osoro2025leo}. However, Kumaran et al. complicate this positive view by shifting attention from individual satellite efficiency to constellation-scale manufacturing impact and orbital occupation. Under their perspective, miniaturization may increase the total number of satellites deployed, and this increase may intensify environmental burden, occupy orbit, and obstruct future launches~\citep{kumaran2024manufacturing}. Kukreja et al. further complicate the access argument by comparing LEO broadband with terrestrial alternatives: even when megaconstellations improve access, their emissions intensity per subscriber remains substantially higher than 4G network systems~\citep{kukreja2025leo}. Therefore, miniaturization is effective when judged by functional efficiency and access, but it is not automatically sustainable when judged by total overall environmental impact.



After analyzing three technological strategies including alternative propellants, reusable rockets and satellite miniaturization, we can find that the disagreement between sources is not simply whether technology is useful, but what level of effect should be evaluated. Dallas, Barato, Hai, and Calabuig support alternative propellants at the level of direct emissions, but Dallas and Calabuig also show that storage, cost, and immature infrastructure limit adoption. Calabuig supports reusable vehicles at the level of per-kilogram launch efficiency, but Wilson's life-cycle perspective and the FAA/McDowell launch cadence data suggest that this efficiency could be weakened by rebound effects. Sandau and Osoro support miniaturization and LEO broadband through access and coverage, but Kumaran and Kukreja shift the discussion toward manufacturing burden, orbital congestion, and emissions intensity per subscriber. By these reasoning, all three technologies are theoretically capable of reducing some environmental impacts, but their real-world effectiveness remains limited by the economic and industrial environment.

The key to addressing this problem is therefore not whether the technology exist, but the conditions under which technology is implemented. Firstly, the market's incentive mechanism was not designed to optimize environmental goals: reducing cost would not necessarily reduce total launches, but may trigger more frequent spaceflight missions and deployments. To be more specific, reusable rockets can escalate spaceflight frequency by reducing its cost, and satellite miniaturization can expand the megaconstellation of broadband networks since it lowers per-satellite manufacturing cost. This scope expansion driven by technological improvements may offset their net environmental benefits, and may even cause reverse effects. Secondly, the existing aerospace industry shows considerable path dependence. For a long time, the rocket industry, including launchpad infrastructure, propellant storage systems, and related supply chains, has been built around traditional fuel systems. This means that although alternative propellants may offer better environmental performance, they may still face considerable limitations in implementation and transformation, since infrastructure-reconstruction costs and industrial inertia are both high. However, even policy itself is not a simple answer. The Outer Space Treaty provides a broad international framework, but Article IX mainly requires states to avoid harmful contamination of outer space and adverse Earth changes from the introduction of extraterrestrial matter; it does not specifically regulate rocket launch emissions or construct a clear emissions standard for commercial megaconstellation launches~\citep[art.~IX]{un1967outerspacetreaty}. This legal gap makes sustainability coordination more complex, since launch providers, satellite operators, and states may all benefit from expansion, while no single international rule directly limits the atmospheric emissions caused by that expansion. By these reasoning, the evidence collectively suggests that the main factor limiting spaceflight sustainability is implementation and transformation, rather than technological capability itself. This also means that, in order to make aerospace activities sustainable, institutional frameworks and policy tools, like limiting launch frequency and guiding industrial transition, are also required, but these policies need stronger international coordination than what current space law provides.



In conclusion, the rapid expansion of LEO broadband satellite constellations has intensified the environmental challenges related to space launches, making sustainability a critical concern for the future of the space industry. While technological advancements, like alternative propellants, reusable rockets and satellite miniaturization, suggest practical approaches to reduce environmental impacts, their effectiveness remains controversial and limited in practice. Alternative propellants provide the most direct way to reduce emissions, but face financial barriers and industrial inertia. Reusable launch vehicles improve per-launch efficiency and also reduce financial cost, but their benefits are limited by the rebound effect that cost reductions may trigger. Similarly, satellite miniaturization improve functional efficiency, but often leads to large-scale deployment that may increase the total environmental burden considerably.


These findings suggest that the primary limitation is not the lack of technological capability, but the implementation of these existing technologies. As long as economic incentives favor expansion and existing industrial systems do not transition, technological improvements alone are not likely to achieve significant reductions in environmental impact. Therefore, ensuring the sustainability of space activities requires not only technological innovation, but also structural changes in political and regulatory frameworks and economic incentives. Finally, sustainable space development depends on promoting the application of technological progress with system-level regulation, rather than relying on methodological advancements alone.

%\subsection{Why Mature Technologies Remain Underutilized}

%\subsection{Barriers to Adoption}

%\subsection{The Case for Regulatory Intervention}

%\section{Conclusion}

\newpage
\bibliographystyle{plainnat}
\bibliography{references}


\end{document}
